The stored Propositions \claimref{std-algebra.polarisation}, \claimref{std-algebra.orthogonal-anticommute}, \claimref{std-algebra.blade-basis} and \claimref{std-algebra.functoriality}, and the existence of $\Cl(V, Q)$ as the quotient of the tensor algebra (\claimref{std-algebra.existence-uniqueness}), are stated over $\R$; their proofs use only field operations and $\tfrac12$, and they are used here verbatim over $F$. Throughout, in a Clifford algebra over $F$, $(u + v)^{2} = Q(u + v)$ expands to $uv + vu = 2B(u, v)$, so two vectors anticommute exactly when $B(u, v) = 0$ (Proposition \claimref{std-algebra.orthogonal-anticommute}, whose proof is the same over any field of characteristic $0$; an ordered field has characteristic $0$ since $m \cdot 1 > 0$ for $m \ge 1$). $B$ and $Q$ determine each other, $Q(v) = B(v, v)$ (Proposition \claimref{std-algebra.polarisation}, again over any such field).

\emph{(i)} Let $a > 0$ and $S = \{x \in F : x \ge 0,\ x^{2} \le a\}$. $S$ contains $m = \min(1, a)$ (if $a \le 1$, $a^{2} \le a$; if $a > 1$, $1 \le a$), and is bounded above by $\max(1, a)$ (if $x > \max(1, a)$ then $x^{2} > x > a$). Let $s = \sup S \ge m > 0$. If $s^{2} < a$, put $h = \min\bigl(1, (a - s^{2})/(2s + 1)\bigr) > 0$; then $(s + h)^{2} = s^{2} + h(2s + h) \le s^{2} + h(2s + 1) \le a$, so $s + h \in S$, contradicting $s = \sup S$. If $s^{2} > a$, put $h = (s^{2} - a)/(2s)$, so $0 < h < s/2$; then $(s - h)^{2} = s^{2} - 2sh + h^{2} > s^{2} - 2sh = a$, and every $x \in S$ has $x \le s - h$ (otherwise $x^{2} > (s - h)^{2} > a$), so $s - h < s$ is an upper bound, contradicting minimality. Hence $s^{2} = a$. The check verifies both inequalities at rational instances, with the hand values $h = \tfrac13$, $(\tfrac43)^{2} = \tfrac{16}{9} \le 2$ and $h = \tfrac1{12}$, $(\tfrac{17}{12})^{2} = \tfrac{289}{144} > 2$ for $a = 2$, and that $2$ is not a rational square ($a^{2} \equiv 2b^{2} \bmod 4$ forces $a$, $b$ even).

\emph{(ii)} Induction on $n = \dim V$. If $B = 0$, every basis consists of mutually anticommuting null vectors. If $B \ne 0$ there is $u$ with $Q(u) \ne 0$: either some basis vector has $Q \ne 0$, or two have $B(x, y) \ne 0$ and then $Q(x + y) = 2B(x, y) \ne 0$. Such $u$ is invertible with $u^{-1} = u/Q(u)$ (since $uu = Q(u)$, a nonzero element of $F$; compare \claimref{std-algebra.inverse-of-vector}), and the rejection $x_{\perp} = x - B(x, u)u^{-1} = \tfrac12(xu - ux)u^{-1}$ satisfies $B(x_{\perp}, u) = B(x, u) - B(x, u)B(u, u)/Q(u) = 0$, so $V = Fu \oplus u^{\perp}$, $x = B(x,u)u^{-1} + x_{\perp}$. By induction $u^{\perp}$ has an anticommuting basis, and with $u$ it is one for $V$, since every vector of $u^{\perp}$ anticommutes with $u$. If $F$ is Euclidean and $Q(u_{i}) \ne 0$, let $\lambda_{i} > 0$ with $\lambda_{i}^{2} = |Q(u_{i})|$; then $Q(u_{i}/\lambda_{i}) = Q(u_{i})/|Q(u_{i})| = \pm 1$, and rescaling keeps anticommutation. The check builds the rejected bases for six Gram matrices (hyperbolic, nondiagonal, degenerate) and verifies in the exact exterior model of the Clifford algebra that $u_{i}u_{j} + u_{j}u_{i} = 2B(u_{i}, u_{j})$ as operators, with $B(u_{i}, u_{j}) = 0$ for $i \ne j$; for $G = A^{T}DA$ with $A$ unitriangular and $D = \operatorname{diag}(4, -9, 0, \tfrac14)$ it recovers $D$ and the normalised squares $1, -1, 0, 1$.

\emph{(iii)} Let $u_{1}, \dots, u_{n}$ be an anticommuting basis, $P$, $N$, $Z$ the spans of the $u_{i}$ with $Q(u_{i}) > 0$, $< 0$, $= 0$, of dimensions $p$, $q$, $r$. For $v = \sum x_{i}u_{i}$, $B(v, u_{j}) = x_{j}Q(u_{j})$, so $v \in \operatorname{rad} V$ exactly when $x_{j} = 0$ for every $j$ with $Q(u_{j}) \ne 0$: $\operatorname{rad} V = Z$ and $r = \dim \operatorname{rad} V$, which does not depend on the basis. On $P$, $Q(\sum x_{i}u_{i}) = \sum x_{i}^{2}Q(u_{i}) > 0$ unless all $x_{i} = 0$, since in an ordered field a sum of positive multiples of squares, not all zero, is positive; so $P$ is a positive definite subspace of dimension $p$. On $N \oplus Z$, $Q \le 0$. If $W$ is positive definite with $\dim W > p$, then $\dim W + \dim(N \oplus Z) > n$; a basis of $W$ together with a basis of $N \oplus Z$ has more than $n$ vectors of $V$, so they are linearly dependent, and a dependence relation gives a nonzero $w \in W \cap (N \oplus Z)$, with $Q(w) > 0$ and $Q(w) \le 0$. So $p$ is the largest dimension of a positive definite subspace, and likewise $q$ for negative definite ones; both are independent of the basis. If $F$ is Euclidean, two spaces with the same $(p, q, r)$ have bases as in (ii) with the same squares; the linear map between them preserves $B$ on basis pairs, hence everywhere, and so preserves $Q$; conversely an isometry carries positive definite subspaces and radicals to the same, so the triples agree. By the universal property, an isometry $V \to F^{p,q,r}$ and its inverse extend to mutually inverse homomorphisms of the Clifford algebras (Proposition \claimref{std-algebra.functoriality}, over $F$), so $\Cl(V, Q) \cong \Cl_{F}(p,q,r)$. The check orthogonalises six Gram matrices from four changes of basis and every ordering of the basis vectors, and finds one triple each, equal to the hand-stated signature, to $r = n - \operatorname{rank} G$, and to the congruence signature of GA.py.

\emph{(iv)} By Theorem \claimref{std-algebra.blade-basis} the $2^{n}$ blades are an $\R$-basis of $\Cl(p,q,r)$, degenerate forms included, and by Proposition \claimref{std-algebra.orthogonal-anticommute} $e_{S}e_{T} = \pm\bigl(\prod_{i \in S \cap T}Q(e_{i})\bigr)e_{S \triangle T}$ with the sign $(-1)^{\#\{(i,j) \in S \times T : i > j\}}$: the structure constants lie in $\{0, 1, -1\} \subset \Q$. They define $\Cl_{\Q}(p,q,r)$, and the $\R$-linear map $\Cl_{\Q}(p,q,r) \otimes_{\Q} \R \to \Cl(p,q,r)$, $e_{S} \otimes \lambda \mapsto \lambda e_{S}$, sends a basis to a basis and respects the products of basis elements, so it is an isomorphism. If $\R$ is complete, it is Euclidean by (i), so by (iii) every real quadratic space has some inertia $(p, q, r)$ and $\Cl(V, Q) \cong \Cl(p,q,r) \cong \Cl_{\Q}(p,q,r) \otimes_{\Q} \R$. Theorem \claimref{conformal.minimal-nondegenerate-completion} embeds $\Cl(p,q,r)$ in $\Cl(p{+}r, q{+}r)$; explicitly, with $r$ added hyperbolic planes spanned by $e'_{k}$, $f'_{k}$ ($Q = 1, -1$), the vectors $e'_{k} + f'_{k}$ are null, anticommute with each other and with the old generators, so the universal property gives a homomorphism, rational on blades, whose images of the $2^{n}$ blades are linearly independent (each has a distinct leading blade). The check verifies that the structure constants of $\Cl(p,q,r)$ lie in $\{0, \pm 1\}$ for four signatures and that the null images realise $\Cl(p,q,r)$ injectively in $\Cl(p{+}r, q{+}r)$ (rank $2^{n}$) for five.

\emph{(v)} Over $\Q$, (ii) without the normalisation and the invariance part of (iii) hold as proved, since $\Q$ is ordered; the isometry criterion of (iii) used the normalisation and fails, as follows. If the line $\Q e$ with $Q(e) = 2$ were isometric to one with $Q(e') = \pm 1$, then $e \mapsto \lambda e'$ with $2 = \pm\lambda^{2}$, so $2$ would be a rational square; writing $\lambda = a/b$ in lowest terms, $a^{2} = 2b^{2}$ forces $a$ even, then $b$ even, a contradiction. In $\Cl(\langle 2 \rangle) = \Q \oplus \Q e$, $x\bar{x} = (a + be)(a - be) = a^{2} - 2b^{2}$, nonzero for $(a, b) \ne 0$ by the same argument, so $x^{-1} = \bar{x}/(x\bar{x})$; in $\Cl_{\Q}(1,0)$, $(1 + e_{1})(1 - e_{1}) = 0$. For the plane, $e_{1}$, $e_{2}$, $e_{1}e_{2}$ anticommute pairwise and $(e_{1}e_{2})^{2} = -e_{1}^{2}e_{2}^{2} = -6$, so for $x = a + be_{1} + ce_{2} + de_{1}e_{2}$ and its Clifford conjugate $\bar{x} = a - be_{1} - ce_{2} - de_{1}e_{2}$ the cross terms cancel and $x\bar{x} = a^{2} - 2b^{2} - 3c^{2} + 6d^{2}$. Suppose it vanishes for rationals not all zero; scale to integers with no common factor. Modulo $3$, $a^{2} + b^{2} \equiv 0$, and as squares are $0$ or $1$ modulo $3$, $3 \mid a$ and $3 \mid b$. Then modulo $9$, $-3(c^{2} - 2d^{2}) \equiv 0$, so $c^{2} + d^{2} \equiv 0 \bmod 3$ and $3 \mid c$, $3 \mid d$, a contradiction. So $x\bar{x} \ne 0$ for $x \ne 0$ and $x^{-1} = \bar{x}/(x\bar{x})$: a division ring. In $\Cl_{\Q}(2,0)$, $(1 + e_{1})(1 - e_{1}) = 0$. Over a Euclidean field, $e_{1}/\sqrt{2}$ and $e_{2}/\sqrt{3}$ are anticommuting vectors of square $1$, so the line and the plane become $\Cl(1,0)$ and $\Cl(2,0)$. The check computes both norms $x\bar{x}$ with GA.py on a grid, confirms that they vanish only at $0$ there and give inverses, runs the two descent lemmas over all residues modulo $4$ and modulo $9$, and exhibits the zero divisors of $\Cl_{\Q}(1,0)$ and $\Cl_{\Q}(2,0)$.

\emph{(vi)} Inspection: (ii) uses only that $Q(u) \ne 0$ can be divided by, and square roots of positive elements for the normalisation; (iii) uses the order (sums of positive multiples of squares are positive) and, for the isometry statement, the normalisation; (iv) uses completeness only to invoke (i).
